\tikzset{
	battery/.style={
		draw=none,
		decorate,
		decoration={
			markings,
			mark=at position 0.5 with {
				% Connectors
				\draw[line] (-\pgfdecoratedcompleteddistance,0) -- (-4,0);
				\draw[line] (4,0) -- (\pgfdecoratedcompleteddistance,0);
				% Battery
				\draw[line] (-4, -20) -- (-4, 20);
				\draw[line] (4, -10) -- (4, 10);
				% Polarity
				\node[above] at (-4, -15) {$+$};
				\node[below] at (4, -15) {$-$};
			}
		}
	}
}

\tikzset{
	switch/.style={
		draw=none,
		decorate,
		decoration={
			markings,
			mark=at position 0.5 with {
				% Connectors
				\draw[line] (-\pgfdecoratedcompleteddistance,0) -- (-15,0);
				\draw[line] (15,0) -- (\pgfdecoratedcompleteddistance,0);
				% Switch
				\draw[line] (-12.5,0) circle (2.5);
				\draw[line] (12.5, 0) circle (2.5);
				\draw[line] (-10,0) -- (10,7);
			}
		}
	}
}

\tikzset{
	spring/.style={
		decorate,
		decoration={
			coil,
			aspect=0.6,
			segment length=7pt,
			amplitude=6pt,
			pre length=5pt,
			post length=5pt
		},
		line width=1.5
	}
}

\tikzset{
	line/.style={
		line width=1.5pt,
	}
}

\tikzset{
	point/.style={
		circle,
		fill=black,
		inner sep=1.5pt
	}
}



%%% THIS IS A HELPER GRID FOR TIKZ PICTURES %%%

%\begin{center}
%	\begin{tikzpicture}
%		
%		% Grid
%		\draw[help lines, step=1] (-2,-2) grid (12,12);
%		
%		% Axes
%		\draw[->, line width=1] (-2,0) -- (12.5,0) node[right] {$x$};
%		\draw[->, line width=1] (0,-2) -- (0,12.5) node[above] {$y$};
%		
%		% Axis labels
%		\foreach \x in {-2,-1,1,2,...,12}
%		\node[below] at (\x,0) {\small $\x$};
%		
%		\foreach \y in {-2,-1,1,2,...,12}
%		\node[left] at (0,\y) {\small $\y$};
%		
%		% Your drawing
%		\draw[spring] (0,0) -- (10,10);
%		
%	\end{tikzpicture}
%\end{center}



%%% ROTATING ARROW HEADS %%%

% To rotate {Stealth} in place use the flex' command:
%\draw[
%line,
%-{Stealth[length=7pt,flex'=0.9]}] 
%(175,186.5) arc (90:270:12.5);



%%% EXAMPLE HATCHING FOR SURFACES %%%
%\begin{scope}
%	\clip (350,100) rectangle (425,275);
%	
%	\foreach \x in {175,183,...,425} {
%		\draw[line width=0.75pt]
%		(\x,100) -- ++(175,+175);
%	}
%\end{scope}